\documentclass{article}
\usepackage[x11names, rgb]{xcolor}
\usepackage[utf8]{inputenc}
\usepackage{tikz}
\usetikzlibrary{arrows,shapes}
\usepackage{amsmath}
%
%
\usepackage[active,tightpage]{preview}
\PreviewEnvironment{tikzpicture}
\setlength\PreviewBorder{5pt}%
%

\definecolor{attackColor}{RGB}{219,6,10}
\definecolor{supportColor}{RGB}{50,144,24}
\definecolor{pro}{RGB}{186, 218, 185}
\definecolor{con}{RGB}{227, 191, 186}

%\tikzstyle{argument}= [rectangle,rounded corners,draw=black, top color=white, bottom color=black!10, thick, text centered]
% for using linebreaks with `\\` we use `align = center` (see: https://tex.stackexchange.com/questions/123671/manual-automatic-line-breaks-and-text-alignment-in-tikz-nodes#124114
%\tikzstyle{argument}= [rectangle,rounded corners,draw=black, top color=white, bottom color=black!10, thick, align=center]
\tikzstyle{argument}= [rectangle,rounded corners,draw=black, top color=white, bottom color=black!10, thick, align=center]
\tikzstyle{argument_pro}= [rectangle,rounded corners,draw=black, top color=white, bottom color=pro, thick, align=center]
\tikzstyle{argument_con}= [rectangle,rounded corners,draw=black, top color=white, bottom color=con, thick, align=center]
%\tikzstyle{thesis}= [rectangle,rounded corners,draw=black, color=white, thick, align=center]
\tikzstyle{thesis}= [rectangle,rounded corners,draw=black, thick, align=center]

\tikzstyle{carneades_statement}= [rectangle, draw=black, thick, align=center]
\tikzstyle{carneades_argument}= [circle,draw=black, thick, align=center]
\tikzstyle{edge_label}= [rectangle, draw=white, thick, align=center]


% fine tuning of edge thickness can be achieved with 'line width=3.8pt'
%\tikzstyle{attack}= [>=latex',  shorten >=1pt, very thick, color=attackColor, auto, text=black]
%\tikzstyle{support}= [>=latex',  shorten >=1pt, very thick, color=supportColor, auto, text=black]

\tikzstyle{attack}= [>=latex',  shorten >=1pt, thick, color=attackColor, auto, text=black]
\tikzstyle{support}= [>=latex',  shorten >=1pt, thick, color=supportColor, auto, text=black]
\tikzstyle{green_edge}= [thick, color=supportColor, auto, text=black]


\begin{document}
\pagestyle{empty}
%
%
%

\enlargethispage{100cm}
% Start of code
% \begin{tikzpicture}[anchor=mid,>=latex',join=bevel,]
\begin{tikzpicture}[>=latex',join=bevel,]
\pgfsetlinewidth{1bp}
%%
\node (C) at (276.0bp,244.0bp) [draw,ellipse,argument] {\textbf{Main claim}\\We should all follow\\ a vegetarian life style.};
  \node (A1) at (188.0bp,156.0bp) [draw,ellipse,argument] {\textbf{Reason $R_1$}\\Most mass meat-farming techniques\\are torture to animals.};
  \node (A2) at (97.0bp,44.0bp) [draw,ellipse,argument] {\textbf{Reason $R_2$}\\Animals feel fear and pain\\like us. Mass meat-farming techniques\\ force animals living in them\\to have an unnaturally\\short, miserable and confined life.};
  % \node (A3) at (280.0bp,44.0bp) [draw,ellipse,argument] {\textbf{Objection $R_3$}\\There is the alternative of\\grass-fed meat, which comes\\from animals that are raised\\in open grass pasture and are\\not confined to cramped\\living spaces.};
  \node (A3) at (280.0bp,44.0bp) [draw,ellipse,argument] {\textbf{Objection $R_3$}\\There are alternative ways to\\breed animals. Grass-fed meat\\comes from animals raised in\\open grass pastures and not\\confined to these cramped\\living spaces.};
  % \node (A4) at (364.0bp,156.0bp) [draw,ellipse,argument] {\textbf{Reason $R_4$}\\The farming of animals produces\\greenhouse gas emissions that\\are causing a dangerous\\climate change.};
  \node (A4) at (364.0bp,156.0bp) [draw,ellipse,argument] {\textbf{Reason $R_4$}\\Animal farming produces more\\greenhouse gas emissions than\\the world's entire transport system,\\and these greenhouse gases are\\causing a dangerous climate change.};
  \draw [->,support] (A1) ..controls (219.01bp,187.3bp) and (235.39bp,203.31bp)  .. (C);
  \draw [->,support] (A2) ..controls (136.62bp,92.887bp) and (153.38bp,113.15bp)  .. (A1);
  \draw [->,attack] (A3) ..controls (232.37bp,101.95bp) and (219.96bp,116.78bp)  .. (A1);
  \draw [->,support] (A4) ..controls (322.69bp,197.37bp) and (312.32bp,207.51bp)  .. (C);
%
\end{tikzpicture}
% End of code

%
\end{document}
%
<<startfigonlysection>>
\begin{tikzpicture}[>=latex,join=bevel,]
\pgfsetlinewidth{1bp}
%%
\node (C) at (276.0bp,244.0bp) [draw,ellipse,argument] {\textbf{Main claim}\\We all should change to\\ a vegetarian life style.};
  \node (A1) at (188.0bp,156.0bp) [draw,ellipse,argument] {\textbf{Reason $R_1$}\\Most mass meat-farming techniques\\are torture to animals.};
  \node (A2) at (97.0bp,44.0bp) [draw,ellipse,argument] {\textbf{Reason $R_2$}\\Animals feel fear and pain\\like us. Mass meat-farming techniques\\ force those animals to have an unnaturally\\short, miserable and confined life.};
  \node (A3) at (280.0bp,44.0bp) [draw,ellipse,argument] {\textbf{Objection $R_3$}\\There is the alternative of\\grass-fed meat, which comes\\from animals that are raised\\in open grass pasture and are\\not confined to cramped\\living spaces.};
  \node (A4) at (364.0bp,156.0bp) [draw,ellipse,argument] {\textbf{Reason $R_4$}\\The farming of animals produces\\greenhouse gas emissions that\\are causing a dangerous\\climate change.};
  \draw [->,support] (A1) ..controls (219.01bp,187.3bp) and (235.39bp,203.31bp)  .. (C);
  \draw [->,support] (A2) ..controls (136.62bp,92.887bp) and (153.38bp,113.15bp)  .. (A1);
  \draw [->,attack] (A3) ..controls (232.37bp,101.95bp) and (219.96bp,116.78bp)  .. (A1);
  \draw [->,support] (A4) ..controls (322.69bp,197.37bp) and (312.32bp,207.51bp)  .. (C);
%
\end{tikzpicture}
<<endfigonlysection>>
